\subsection{Version 2: contemporary solutions and attribution}
\label{sec:v2-changes}
The following changes were made on 19 September after the editor
identified further contemporary solutions. They compare the complete
pre-v2 paper at local commit \artifact{5fcdfdf} with version~2.
The preserved PDF is
\artifact{paper/versions/kourovka-experiment-2026-09-17-v1.pdf}.
Source hashes and literal excerpts are bound in
\artifact{paper/reviews/v2-2026-09-19/}. Each pair identifies the old
statement and its replacement or the additional qualification.

\paragraph{The task input}

\noindent\textbf{Before.}
\begin{quote}\small
48 hours to investigate previously unsolved problems in the Kourovka
Notebook
\end{quote}

\noindent\textbf{Version 2.}
\begin{quote}\small
48 hours to investigate questions listed as open in the supplied Kourovka
Notebook
\end{quote}

State what the supplied edition listed; later attribution does not establish that every question was still open at the time of the run.

\paragraph{The abstract now reports the comparison}

\noindent\textbf{Before.}
\begin{quote}\small
finite certificates, with their hypotheses and imported ingredients.
\end{quote}

\noindent\textbf{Version 2.}
\begin{quote}\small
We compare eleven entries with contemporary solutions, identifying shared
arguments, different constructions and stronger alternatives.
\end{quote}

Add the principal outcome of the second version immediately after the mathematical summary.

\paragraph{Coverage and novelty}

\noindent\textbf{Before.}
\begin{quote}\small
In particular, we acknowledge
contemporaneous work on 21.106, and retain known results and deductions
under their prior-work status.
\end{quote}

\noindent\textbf{Version 2.}
\begin{quote}\small
We assign no discovery priority from the available chronology, and do not
infer that the other 35 entries are new.
\end{quote}

Replace the single-overlap discussion with an eleven-entry comparison while retaining 46 as a historical classification.

\paragraph{16.9: prior qualitative computability}

\noindent\textbf{Before.}
\begin{quote}\small
novelty is a separate question.
\end{quote}

\noindent\textbf{Version 2.}
\begin{quote}\small
the qualitative computability assertion
in 16.9 is now explicitly identified as a consequence of older work.
\end{quote}

The new attribution note credits Dahmani-Guirardel for the decision procedure and Brandenbursky-Gal-Kedra-Marcinkowski for the interval recurrence, following the contemporary paper's source discussion.

\paragraph{Specialist confirmation has a specific scope}

\noindent\textbf{Before.}
\begin{quote}\small
Its independence concerns the reviewing model and the separately written
checks; it is not a proof of discovery independence, human specialist
endorsement or formal verification.
\end{quote}

\noindent\textbf{Version 2.}
\begin{quote}\small
He reported that group theorists had confirmed alternative
solutions for 21.68, 16.28(a) and 15.89; many other drafts remained
unverified.
\end{quote}

Add the editor's report about particular alternatives without changing the scope of Claude's earlier review.


\paragraph{Individual attribution and evidence.}
Eleven new notes accompany the affected proofs. They credit the papers
listed in Appendix~\ref{app:contemporary-comparison}, including both
Rizzoli and Jayadevan for 21.68 and the retained Zhang--Li work for
21.106. The inventory identifies every overlap. The new appendix gives
the comparison, chronology and checks; the Fable appendix now explicitly
retains the historical meaning of our answer count in that comparison.
The original proof expositions, deadline ledger and four exchanged
review documents remain preserved. The v2 comparison and algebraic
identities are later work, separate from the deadline results and
Claude's original review. The raw editor email is not reproduced.
